\documentclass[11pt]{article}
\usepackage{fontspec,amsmath,amssymb,amsthm,hyperref}
\usepackage[margin=1in]{geometry}
\setlength{\emergencystretch}{3em}
\title{Arbitrary-friction contraction directly from the potential}
\author{LangevinLean mathematical exposition}
\date{}
\begin{document}\maketitle
\noindent Source module: \texttt{KineticLangevin/GeneralContraction.lean}. Compile this proof document with XeLaTeX or LuaLaTeX. The accompanying HTML pairs the mathematical argument with the exact Lean statements and records the scope of the formalization.


The paper's full shifted contraction can be derived from differentiability, strong convexity and a Lipschitz gradient alone. This proof centers velocity at \(v+\gamma x/2\), uses the strong-smooth sector inequality, and completes a scalar square. It removes any need to assume that the Hessian is continuous when constructing a mean Hessian, while recovering the exact maximum of the two endpoint curvature exponents.


\section{Center the phase energy and its derivative}


The concrete force-difference derivative is \(2\langle v+(\gamma-a)x,-av-F\rangle+2\langle-v-ax,(\gamma-a)v+F\rangle\). Put \(y=v+\gamma x/2\), \(d=\gamma-2a\), and \(c=((\gamma-a)^2+a^2)/2\). Expanding the inner products gives the two exact identities

\[\mathcal E=2\|y\|^2+\tfrac12d^2\|x\|^2,\qquad \mathcal E'=-\gamma\mathcal E+4\langle y,cx-F\rangle.\]

The Lean proof rewrites squared norms as self-inner-products, distributes inner products over sums, differences and scalar multiplication, uses symmetry, and finishes with \texttt{ring}. These identities expose a scalar residual rather than a matrix diagonalization.



\noindent Lean declarations: \texttt{generalForceEnergyDerivative}, \texttt{transformEnergy\_centered}, \texttt{generalForceEnergyDerivative\_centered}.

\section{Bound the residual around an arbitrary scalar curvature}


Strong-smooth interpolation places \(F\) in the ball centered at \(\mu x\), where \(\mu=(m+M)/2\), with radius \(R\|x\|\), where \(R=(M-m)/2\). Decompose \(cx-F=(c-\mu)x-(F-\mu x)\). The triangle inequality gives \(\|cx-F\|\le(|c-\mu|+R)\|x\|\). The coefficient is exactly \(\max\{c-m,M-c\}\): Lean splits on \(c\ge\mu\), simplifies the absolute value and the corresponding maximum, and checks the remaining arithmetic. No gradient or Hessian operator is invented in this argument.



\noindent Lean declarations: \texttt{centered\_force\_residual}.

\section{Complete the centered square to obtain the exact exponent}


Let \(K=\max\{c-m,M-c\}\ge0\). Cauchy–Schwarz and the residual bound give \(\langle y,cx-F\rangle\le K\|y\|\|x\|\). The nonnegative square \((2\|y\|-d\|x\|)^2\) implies \(4d\|y\|\|x\|\le2\mathcal E\). Multiplying by \(K\) and inserting the exact energy-derivative identity yields \(d\mathcal E'\le(-\gamma d+2K)\mathcal E\). Finally \(K=\max\{a^2-m,M-(\gamma-a)^2\}+\gamma d/2\), so the \(\gamma d\) terms cancel. Lean records the multiplied identities explicitly before \texttt{nlinarith}, then divides by the positive \(d\). This proves the precise spectral differential estimate even in a friction regime with a positive growth exponent.



\noindent Lean declarations: \texttt{general\_force\_dissipativity}.

\section{Supply the sector inequality from the actual potential}


For \(x-y\) and \(F=\nabla f(x)-\nabla f(y)\), the project's proved \texttt{strong\_smooth\_gradient\_interpolation} gives the strong-smooth interpolation inequality from \texttt{Differentiable}, \texttt{StrongConvexOn} and \texttt{LipschitzWith}. Its centered-ball consequence is exactly the residual theorem's premise. The proof composes these already proved results, yielding arbitrary-shift dissipativity for the actual nonlinear gradient. It assumes neither a Hessian representation nor the desired contraction conclusion.



\noindent Lean declarations: \texttt{potential\_general\_force\_dissipativity}.

\section{Integrate the right-derivative inequality}


The two kinetic difference equations imply the transformed-coordinate right derivatives \(-av-F\) and \((\gamma-a)v+F\). Their squared-norm derivative rules add to \texttt{generalForceEnergyDerivative}. For paths on \([0,T]\), define their transformed energy \(E(t)\). Continuity follows from continuity of the paths, while the preceding nonlinear result gives \(E'(t)\le2kE(t)\), with \(k=\max\{a^2-m,M-(\gamma-a)^2\}/(\gamma-2a)\). The right-slope form of Grönwall then proves \(E(t)\le E(0)e^{2kt}\), including the initial endpoint. The hypotheses are the actual force-difference differential equations, not an exponential bound supplied by the caller.



\noindent Lean declarations: \texttt{hasDerivWithinAt\_generalForceEnergy}, \texttt{potential\_general\_energy\_decay\_right}.

\end{document}
